\scanpage{258}

\diagram{\begin{tikzcd}[ampersand replacement=\&,column sep=large,row sep=large]
\& \widetilde Z\times\widetilde Z \arrow[r] \arrow[d,"{(\widetilde\lambda,\widetilde\mu)\mapsto\widetilde\lambda+\widetilde\mu}"] \& T_G\arrow[r]\arrow[d,"\delta"] \& 1\\
(123) \& \widetilde Z\arrow[r,"c"] \& Z. \&
\end{tikzcd}}

Reste à décrire, \unsure{à} part \deleted{\unsure{l'incl|l'action}} (115) lui-$=$, et l'action de $T_G$ sur $H$ \inserted{(donnée c))}, i.e. de $T_0\xrightarrow{\sim}\widetilde Z$ et de $T_1\simeq\widetilde Z$ sur $H$, l'inclusion de $Z$ dans $T_G$ (donnée d) de (116')). \deleted{Je dis que $\unsure{c}$} \inserted{Cette inclusion} peut se décrire à l'aide de l'hom

\diagram{\begin{tikzcd}[ampersand replacement=\&]
(12{\deleted{\unsure{6}}4})\qquad c':Z \arrow[rr] \arrow[dr,swap,"{\deleted{\unsure{\mathrm{incl}.}}}","{\diagup}" description] \& \& \widetilde Z,\\
\& {\deleted{\unsure{T_G}}}\arrow[ur,"{\deleted{\unsure{\widetilde\delta}}}","{\diagup}" description] \&
\end{tikzcd}}

défini comme composé

\[Z\overset{\mathrm{incl}}{\unsure{-|\longrightarrow}}\deleted{\unsure{T_G}}G\xrightarrow{\widetilde\delta}\widetilde Z\quad\text{ou}\quad Z\overset{\mathrm{incl}}{\unsure{-|\longrightarrow}}T_G\xrightarrow{\widetilde\delta|T_G}\widetilde Z,\tag*{(12\unsure{5|4})}\]

qui satisfait la relation

\diagram{\begin{tikzcd}[ampersand replacement=\&,column sep=large]
(12\unsure{6|5})\quad Z \arrow[r,"c'"] \arrow[rr,bend right=24,"{\lambda\mapsto\lambda^2}" description] \& \widetilde Z\arrow[r,"c"] \& Z\qquad \begin{aligned}&\text{i.e.}\\&c(c'(\lambda))=\lambda^2\\&\forall\lambda\in Z.\end{aligned}
\end{tikzcd}}

Ainsi, quand on écrit

\[T_G=\deleted{Z}T_H\times\deleted{\widetilde Z}T_0\tag*{(12\unsure{6|5|7})}\]

la donnée de l'inclusion $Z\subset T_G$ revient à celle d'un hom

\[Z\hookrightarrow\underset{Z}{\underline{T_H}}\times\underset{\widetilde Z}{\underline{T_0}},\tag*{(12\unsure{7|6})}\]

donné par la compatibilité

% archive sheet number: 257
